\chapter{SYSTEM ANALYSIS}
\par
\hspace{0.5cm}
System analysis is the process of structurally examining the proposed AI-Powered Code Document Generator system. The main purpose of the analysis is to comprehend the system’s functionality, processing flow and interaction in the creation of accurate and updated documentation from software source code automatically.

\par The analysis includes an analysis based on how the user uses the whole system, such as a source code directory or a version control system can be provided as input. This includes parsing of the source code, feature extraction and then finally document generation with the help of a LLM. On this case, it is evident that the analysis leads to the production of documents in an efficient way and can be easily maintained and scaled to usable software development purposes. It also ensures an appropriate design of the system that will specify functional and non-functional requirements efficiently and will ensure that the proposed system meets all the user requirements while maintaining accuracy and performance. 

\section{Requirements Analysis}
\hspace{0.5cm}
The requirements analysis gives a broad functional and user requirements for documentation using a  coding system. It is a requirement that this system will be capable of analyzing the software source code and producing documentation that is easily understood by human readers. All the requirements are specified in the most clear way to be precise, measurable and attainable by the system. The main requirements are that the application should be able to parse source code files, extract structural components of the program, integrate the language model and provide a user interface to display the documentation. Also, the system should be able to support many programming languages and optimally process different project structure types. These requirements are mandatory to design a reliable, efficient and intuitive system which can generate a large-scale documentation of source code without causing a failure of user expectations and maintaining a level of accuracy and performance. 
\section{Existing System}
\hspace{0.5cm}
As of now, software documentation management systems are either based on manual documentation processes or use simple automated documentation systems. In the manual process, programmers create documentation by writing comments, README files and design documents manually. While this approach allows for customization, it requires a significant amount of time and heavily depends on the individual programmer’s effort and expertise. Automated documentation systems are also employed to reduce the time required for writing software documentation. These systems generally generate documentation automatically, either using code comments or predefined templates. However, most of these automated systems rely on superficial methods and lack the capability to understand the deeper structural and semantic aspects of the code, limiting their effectiveness in producing thorough and accurate documentation.


\section{Limitations of Existing System}
\hspace{0.5cm}
Though there is access to documentation technology, the existing systems have the following drawbacks:

\begin{itemize}
\item Documentation is usually left up to developers to manually update, especially during fast development cycles. Such situations can cause confusion for new developers and hinder maintenance efforts.
\item The tools themselves may not fully understand the logic, data flow or functionality of the code, limiting the generation of insightful and meaningful documentation.
\item Some documentation tools struggle to capture the relationships and connections between classes, modules and functions, making it difficult to map out the overall program structure.
\item Many current documentation methods support only a single programming language, reducing flexibility and applicability in multi-language projects. In such cases, documentation may become fragmented or inconsistent.
\item Automatically generated documentation often lacks context about the code, leaving the purpose, function and interactions of components unclear, which diminishes its usefulness as a reference for developers.
\end{itemize}

\section{Proposed System}
\hspace{0.5cm}
The proposed system is based on an AI-Powered Code Document Generator that is meant to produce accurate, structured and up-to-date documents automatically, based on the software source code. Differently from other systems, this system combines advanced code analysis with the use of a LLM to understand the structure and semantics of the code. The system accepts input in the form of a local source code folder or a version-controlled repository. Next, the system can use the extracted features to generate meaningful documents such as file descriptions, architectural descriptions or project descriptions. 

\section{Advantages of Proposed System}
\hspace{0.5cm}
The proposed system offers the following advantages:

\begin{itemize}
\item Minimal manual work is required for documentation generation, as the process is fully automated. This allows developers to focus more on coding rather than spending time writing extensive documentation.
\item Precise understanding of the program’s code structure is achieved through AST-based analysis, enabling accurate extraction of functions, classes and dependencies. The system provides detailed insights into both the structure and functionality of the code.
\item Documentation is human-friendly and context-aware, generated with consideration of the purpose and interactions of specific code elements. This helps developers read and understand the system structure more quickly and efficiently.
\item Multi-language support increases versatility, especially for teams working in polyglot environments, eliminating the need for separate documentation tools for different programming languages.
\item Efficient handling of large and complex code repositories ensures that the system can provide comprehensive documentation across entire codebases, aiding in the maintenance and navigation of large projects.
\item Enhanced understandability and maintainability of source code improve collaboration within teams, ensuring that documentation supports long-term project sustainability and effective knowledge sharing.
\end{itemize}
\section{Functional Requirements}
\hspace{0.5cm}
This section describes the core functionalities that the system must provide to meet user requirements. Functional requirements define the specific operations and behaviors of the system when analyzing source code and generating documentation. They outline what the system should do and serve as a foundation for designing and developing its various modules. 

\subsection{User Interface}
\hspace{0.5cm}
This section offers a summary of the user-system interaction layer. Through a user system interface, the user shall specify the location of local copies of source code files or version control repositories for generating the documentation. This includes specification of the location of local source code files or version-controlled repositories for documentation generation and the interface will ensure easy interaction with no required technical knowledge or skills. 

\subsection{Source Code Input}
\hspace{0.5cm}
This section outlines the supported input formats. The system should accept software projects provided as local directories or version-controlled repositories. It must handle various repository formats and support multiple file types within each project. Additionally, the system should validate all inputs to ensure correctness before initiating the analysis process.

\subsection{Source Code Parsing}
\hspace{0.5cm}
This subsection shows the method that is used for internal source code analysis. To ensure correct understanding of hierarchical code, source code files are parsed through AST generation. Correct parsing results are ensured to identify all the relevant components to work on further correctly. 

\subsection{Feature Extraction}
\hspace{0.5cm}
The following section is a description of how meaningful features are extracted from the code. The system parses the source code for structural elements like modules, classes, methods and functions, imports or dependencies. The code is parsed to extract structured elements such as modules, classes and functions, along with imports and dependencies. 

\subsection{LLM-Based Documentation Generation}
\hspace{0.5cm}
This subsection describes documentation creation processes. For documentation creation, a dedicated LLM is trained to convert the previously extracted code properties into human-readable documentation. To give the most appropriate human-readable documentation, a specialized LLM considers the structure and the context of the generated output during the documentation creation procedure. 

\subsection{Documentation Output}
\hspace{0.5cm}
The following section describes the various document formats produced by this system will hence be described in the following paragraphs. These are a README file, a description of the architecture, a release of the latest changes and a description at the file level. This means that the output shall be given in such a way that it is concise and clear, making it easier to analyze the document produced. 

These are the generated files:
\begin{itemize}
\item \textbf{README.md} – Provides a detailed description of the project, including its name, programming language(s), total functions and classes, key imports, overall project description, file structure and instructions for running the code.
\item \textbf{architecture\_overview.md} – Offers an overview of the system architecture, including per-file responsibilities, module dependencies, inter-module interactions and identification of critical or unused functions.
\item \textbf{recent\_changes.md} – Lists recent commits, including authors and commit messages, summarizing updates and modifications in the project.
\item \textbf{files/*.md} – Contains documentation for each source code file, detailing namespaces, classes, functions, member variables, imports, internal and cross-file relationships, entry points and potential issues such as recursion or unused functions.
\end{itemize}


\section{Non-Functional Requirements}
\hspace{0.5cm}
The non-functional requirements specify quality attributes for the system. They ensure that the systems output is accurate, consistent and that it functions effectively and remains user-friendly. The non-functional requirements specify the quality attributes for the system to ensure that it is effective, user-friendly and produces the right output and that it emphasizes on scalable, usability and reliability that underpins implementing a robust and practical AI-powered documentation solution. 
\subsection{Scalability}
\hspace{0.5cm}This system must be able to work with both small projects in source code files and large projects in multi-module repositories. The system’s efficiency and uniformity are not compromised with the increasing volume of source codes. The modular design enables multiple files to be simultaneously processed for efficiency and to maintain consistency in performance for large-scale applications. 

\subsection{Usability}
\hspace{0.5cm}The system must provide an easy, straightforward interface to allow users either to upload their source code or provide Uniform Resource Locator (URL) of repositories. The interface should be intuitive and self-explanatory, such that the user is able to generate the documentation without the need for any complex configuration, thereby minimizing any learning requirements for new users.

\subsection{Reliability}
\hspace{0.5cm}Documentation that is generated should be accurate and well structured. It should be able to cope with such issues as inappropriate or unsupported files without breaking the workflow. This requires an efficient error-handling mechanism that prevents the system from crashing and provides a consistent, stable and reliable performance. 

\subsection{Maintainability}
\hspace{0.5cm}The system’s architecture should allow for its customization and upgrading when changes in programming languages, libraries or documentation norms occur. It is done by using a modular architecture, well-structured code and configurable components that will help the developers customize or upgrade the features with minimum effort, guaranteeing the long-term usability of the documentation system. 

\section{Feasibility Study}
\hspace{0.5cm}
This section discusses the feasibility of the proposed system from technical, operational and economic perspectives. The system is considered viable because it uses open-source tools and locally hosted language models, eliminating licensing costs and ensuring practical, cost-effective implementation.

\subsection{Technical Feasibility}
\hspace{0.5cm}
This subsection evaluates the technical feasibility of implementing the proposed system. It relies on widely used and mature technologies, including Python, AST-based parsers and LLMs, all supported by active and robust communities. The system’s moderate computational requirements allow it to be developed and executed efficiently in standard computing environments.
\subsection{Operational Feasibility}
\hspace{0.5cm}
This subsection analyzes the usability of the proposed system. It requires minimal user input and little technical knowledge, as users only need to specify the location of the source code file, while the system handles all processing automatically.

\subsection{Economical Feasibility}
\hspace{0.5cm}
This subsection analyzes cost-related aspects. The system uses open-source software and libraries, avoiding any licensing fees and operates on standard hardware, eliminating additional infrastructure costs. This makes it affordable, cost-effective and accessible for both individual users and institutions.
\newline
\par The analysis of the proposed AI-Powered Code Documentation Generator was carried out by thoroughly examining the system requirements, identifying the limitations of existing documentation solutions and evaluating the functional features of the proposed system. This analysis helps in understanding how the system is expected to behave under different scenarios and ensures that all user requirements are adequately addressed. Furthermore, it highlights the gaps in current approaches and justifies the need for an intelligent, automated solution. Overall, this comprehensive analysis provides clear insight into the system’s behavior and establishes a strong foundation for the system design presented in the following chapter.

% \subsection{Error Handling}
% Invalid inputs, unsupported files, or analysis failures shall be
% detected and communicated to the user through appropriate error
% messages.

% \subsection{Data Security}
% The system shall process source code locally and shall not permanently
% store or transmit user code without authorization, ensuring data
% privacy and security.

% \section{Feasibility Study}
% \hspace{0.5cm}A feasibility study evaluates the practicality of implementing the
% proposed system in terms of technical, operational, and economic
% factors. This assessment ensures that the system can be developed using
% available resources within reasonable constraints.

% \subsection{Technical Feasibility}
% \hspace{0.5cm}The system utilizes widely supported technologies such as Python,
% AST-based parsers, vector databases, and large language models. The
% computational requirements are moderate and can be satisfied using
% standard personal computing environments, making the system technically
% feasible.

% \subsection{Operational Feasibility}
% \hspace{0.5cm}The system is designed to require minimal user interaction and technical
% expertise. Users only need to provide a source code location to obtain
% generated documentation, making the system operationally feasible for
% students, developers, and academic users.

% \subsection{Economical Feasibility}
% \hspace{0.5cm}The system is developed using open-source tools and libraries. No
% specialized hardware or licensed software is required, making the
% proposed solution cost-effective and economically feasible.
